\chapter{Apathy}

\section*{Definition}
Apathy is a persistent reduction in motivation, interest, emotional responsiveness, or goal-directed activity. It is not the same as laziness. Apathy and depression can overlap, and depression may occur without obvious sadness, so they cannot be separated by one symptom alone.

\section*{When to See a Doctor}

\begin{itemize}
 \item Apathy appears suddenly or is accompanied by confusion, cognitive decline, new motor symptoms, or personality change
 \item It significantly impairs work, relationships, or self-care
 \item There is concern for an underlying neurological or medical disorder
\end{itemize}

\section*{How It Develops}
Apathy may reflect disruption in interacting brain networks involved in motivation, reward, attention, and executive function, including frontal and basal-ganglia circuits. Depression, medication effects, sleep problems, illness, and social or environmental factors can produce similar behavior. No single neurotransmitter test or brain scan diagnoses apathy.

\section*{Causes}
\begin{itemize}
 \item \textbf{Neurodegenerative:} Alzheimer disease, frontotemporal dementia, Parkinson disease, progressive supranuclear palsy, Huntington disease
 \item \textbf{Psychiatric:} Schizophrenia (negative symptoms), major depressive disorder, chronic stress
 \item \textbf{Neurological:} Traumatic brain injury, stroke (especially basal ganglia or frontal lobe), multiple sclerosis, HIV encephalopathy
 \item \textbf{Endocrine/metabolic:} Hypothyroidism, vitamin B12 deficiency, anemia, nutritional deficiency, or other systemic illness; low testosterone is considered only when symptoms and repeated testing support it
 \item \textbf{Medication-induced:} Antipsychotics (D2 blockade), SSRIs, benzodiazepines, beta-blockers
\end{itemize}

\section*{Related Symptoms}
\begin{itemize}
 \item Anhedonia, social withdrawal, abulia, flat affect, psychomotor retardation, fatigue, cognitive impairment
\end{itemize}
\textbf{Important:} Apathy and depression frequently co-occur but require different treatment approaches; screening for depression is essential.

\section*{Medical Evaluation}
\begin{itemize}
 \item Clinical history differentiating apathy from depression, anhedonia, fatigue, sleepiness, medication effects, and cognitive impairment, ideally with information from someone who knows the person well
 \item Structured scales: Apathy Evaluation Scale (AES), Lille Apathy Rating Scale (LARS)
 \item Cognitive screening (MoCA, MMSE) and neurological exam to identify underlying disorders
 \item Targeted tests such as blood count, TSH, B12, metabolic studies, or infection testing when indicated; testosterone testing is not routine and should be interpreted with symptoms and repeat morning measurements
 \item Neuroimaging if focal neurological signs or rapid progression
\end{itemize}

\section*{Treatment Options}
\begin{itemize}
 \item Treat the underlying condition, such as depression, sleep disorder, thyroid disease, medication effect, or Parkinson disease. Testosterone replacement is considered only for confirmed deficiency with compatible symptoms and appropriate monitoring.
 \item Non-pharmacological: behavioral activation, cognitive rehabilitation, environmental enrichment
 \item Psychostimulants or dopaminergic medicines may be considered off-label by specialists in selected cases, but evidence is limited and risks include insomnia, appetite loss, cardiovascular effects, impulsivity, and worsening psychosis.
 \item For dementia-related apathy, activity-based and psychosocial approaches are preferred; cognitive medicines should be used for an indicated dementia diagnosis, not routinely for apathy alone.
\end{itemize}

\section*{Self-Care and Home Management}
\begin{itemize}
 \item Establish structured daily routines with specific, attainable goals
 \item Use activity scheduling and behavioral activation techniques
 \item Engage in social activities and exercise, starting with small commitments
 \item Avoid over-sedating substances and review medications with your doctor
\end{itemize}

\section*{References}
\begin{thebibliography}{99}
\bibitem{apathy:ref1} Levy R, Dubois B. ``Apathy and the functional anatomy of the prefrontal cortex--basal ganglia circuits.'' \textit{Cereb Cortex}. 2006;16(7):916-928.
\bibitem{apathy:ref2} Husain M, Roiser JP. ``Neuroscience of apathy and anhedonia.'' \textit{Nat Rev Neurosci}. 2018;19(8):470-484.
\bibitem{apathy:ref3} Lanctôt KL, Agüera-Ortiz L, et al. ``Apathy associated with neurocognitive disorders: recent progress and future directions.'' \textit{Alzheimers Dement}. 2017;13(1):84-100.
\bibitem{apathy:ref4} Miller BL, Boeve BF, eds. \textit{The Behavioral Neurology of Dementia}, 2nd ed. Cambridge University Press, 2016.
\bibitem{apathy:ref5} UpToDate. ``Apathy in dementia.'' Wolters Kluwer, 2023.
\bibitem{apathy:ref6} National Institute for Health and Care Excellence. Dementia: assessment, management and support for people living with dementia and their carers (NG97). Updated 2025.
\end{thebibliography}
